\chapter{Introduction}
\label{ch:introduction}

Cardiovascular diseases remain the leading contributor to global mortality and long-term morbidity, accounting for nearly eighteen million deaths annually according to the World Health Organization. Among the acute manifestations of cardiovascular pathology, Acute Myocardial Infarction (AMI)—colloquially designated as a heart attack—represents a time-critical medical emergency. In clinical practice, an AMI occurs when myocardial tissue suffers prolonged ischemia and necrosis due to an acute imbalance between coronary blood supply and myocardial oxygen demand, predominantly initiated by atherosclerotic plaque rupture and subsequent thrombotic coronary arterial occlusion \citep{thygesen2018fourth}. The survival rate, post-infarction functional status, and long-term recovery of a patient admitted with acute coronary syndromes hinge decisively upon rapid clinical assessment, accurate risk stratification, and the timely administration of guideline-directed interventions such as primary percutaneous coronary intervention or thrombolysis.

Over the past decade, the rapid digitization of hospital intensive care units (ICUs) and the ubiquitous adoption of Electronic Health Record (EHR) systems have produced massive streams of continuous, high-frequency physiological data. Standard bedside patient monitoring systems continuously measure fundamental vital signs, including Heart Rate (HR), Systolic Blood Pressure (SBP), Diastolic Blood Pressure (DBP), Mean Arterial Blood Pressure (MAP), and Peripheral Oxygen Saturation ($\text{SpO}_2$). Because these physiological parameters reflect real-time autonomic balance, hemodynamic stability, and peripheral tissue perfusion, they contain rich latent signatures of cardiovascular stress and decompensation. Modern machine learning techniques offer an unprecedented capability to analyze these high-dimensional physiological time-series, uncovering subtle clinical trajectories that may precede catastrophic cardiovascular decline.

However, translating machine learning research into clinical reality encounters a formidable obstacle: the sensitive nature of patient health records and the strict regulatory frameworks governing biomedical data. Legislation such as the Health Insurance Portability and Accountability Act (HIPAA) in the United States and the General Data Protection Regulation (GDPR) in the European Union strictly restrict the unauthorized aggregation, cross-border sharing, and centralized storage of identifiable clinical records. Healthcare institutions are legally and ethically obligated to safeguard patient privacy, creating isolated ``data silos.'' When individual hospitals attempt to train machine learning models strictly on their internal patient cohorts, they frequently suffer from small sample sizes, demographic bias, and statistical overfitting. A machine learning model trained on a small, localized patient pool typically fails to generalize when deployed at external hospitals with different patient populations and measurement distributions.

To resolve this fundamental tension between data privacy and model generalizability, Federated Learning (FL) has emerged as an innovative distributed machine learning paradigm \citep{mcmahan2017communication, rieke2020future}. Rather than transmitting raw patient records to a centralized computing cluster, Federated Learning brings the model training to the data. Participating healthcare centers retain complete ownership and physical custody of their sensitive clinical databases. Each hospital locally optimizes a private copy of the model using its internal patient records and transmits only anonymized mathematical parameters (such as weight tensors and bias vectors) to a coordinating aggregation server. The central server synthesizes these decentralized updates using aggregation algorithms such as Federated Averaging (FedAvg) and broadcasts an improved global model back to the institutions.

This thesis presents a comprehensive, reproducible, and privacy-preserving framework for cardiovascular admission classification using continuous intensive care vital signs extracted from the MIMIC-IV clinical database (version 3.1). By transforming continuous physiological time-series into fixed-length statistical representations and evaluating both centralized and decentralized architectures, this investigation demonstrates how collaborative multi-hospital machine learning can be achieved without compromising patient confidentiality.

\section{Motivation and Objectives}
\label{sec:motivation_objectives}

The core motivation of this investigation arises from two interrelated challenges currently confronting clinical healthcare informatics:

\begin{enumerate}
    \item \textbf{The Data Silo Bottleneck in Healthcare AI:} While deep and complex machine learning models require diverse patient records to learn robust diagnostic decision boundaries, real-world hospitals cannot pool their raw patient databases due to privacy legislation, institutional competition, and patient consent restrictions. Small hospitals with limited patient intakes are particularly disadvantaged; their localized models often experience severe performance collapse when evaluated on broader patient demographics. There is an urgent practical demand for collaborative frameworks that permit inter-institutional knowledge sharing while guaranteeing strict data residency.
    
    \item \textbf{Early Non-Invasive Hemodynamic Risk Stratification:} In clinical emergency departments and intensive care units, advanced laboratory diagnostic markers such as cardiac troponin often require substantial processing latency (ranging from forty-five minutes to several hours), and invasive angiographic imaging is not immediately accessible in resource-constrained community clinics. Conversely, routine vital signs (heart rate, blood pressure, oxygen saturation) are continuously, non-invasively, and immediately recorded from the exact moment of admission. Demonstrating what predictive capacity resides strictly within the first 24 hours of non-invasive vital signs provides valuable clinical insight into early automated triage.
\end{enumerate}

To directly address these challenges, this thesis defines the following concrete academic and engineering objectives:

\begin{itemize}
    \item \textbf{Cohort Extraction and Clinical Preprocessing:} Systematically explore the MIMIC-IV v3.1 database to extract an adjudicated cohort of cardiovascular admissions identified via International Classification of Diseases (ICD-9 and ICD-10) diagnosis codes, isolating high-frequency timestamped chart events for heart rate, blood pressure, and pulse oximetry while excluding troponin laboratory values from the predictive feature space.
    
    \item \textbf{Leakage-Free Feature Engineering (One Row Per Patient):} Design a rigorous, reproducible preprocessing pipeline that restricts observations strictly to the first 24 hours of hospital vital monitoring (preventing lookahead bias) and transforms multi-rate, irregular physiological time-series into fixed-length tabular vectors summarizing hemodynamic distribution, signal variability, linear rate of change (trend), and clinically established composite ratios such as the Shock Index and Pulse Pressure.
    
    \item \textbf{Centralized Machine Learning Benchmarking:} Implement, tune, and evaluate six foundational machine learning algorithms—Logistic Regression, Random Forest, Decision Tree, Support Vector Machine (SVM), K-Nearest Neighbours (KNN), and Gradient Boosting—using five-fold stratified cross-validation on training data to establish an empirical performance ceiling and rigorously analyze why certain algorithmic families succeed or fail on small-sample medical data.
    
    \item \textbf{Decentralized Multi-Hospital Simulation:} Formulate an authentic distributed simulation comprising three distinct hospital clients (Hospital A: Urban Medical Center, Hospital B: University Hospital, and Hospital C: Regional Health Center), assigning patient records strictly at the subject level so that no patient's data crosses institutional firewalls.
    
    \item \textbf{Mathematical Formulation and Execution of FedAvg:} Implement a PyTorch-based Federated Learning architecture executing Federated Averaging on compatible model parameters across multiple communication rounds, evaluating convergence behavior under both Independent and Identically Distributed (IID) and non-homogeneous (Non-IID) institutional patient prevalence distributions.
    
    \item \textbf{Empirical Baseline Comparison and Isolated Hospital Analysis:} Train independent baseline models at each hospital using only its internal records and quantify the exact performance gap between isolated hospital training, pooled centralized learning, and collaborative federated learning on an identical, strictly held-out global test cohort.
    
    \item \textbf{Interactive Web Application Deployment:} Build and deploy a complete, full-stack research web application (utilizing Python Flask, modern responsive interfaces, and REST API endpoints) that encapsulates the exact trained scaling parameters, feature extractors, and federated models to deliver real-time clinical probability assessment, hemodynamic risk visualization, and federated topology inspection.
\end{itemize}

\section{Report Organisation}
\label{sec:report_organisation}

The remainder of this thesis is systematically organized into the following chapters:

\begin{itemize}
    \item \textbf{Chapter 2: Related Works} presents an exhaustive review of published literature concerning clinical cardiovascular risk assessment, traditional machine learning applications in intensive care units, healthcare data privacy legislation, and the theoretical emergence of federated learning. It includes a comprehensive comparative summary table of landmark studies, outlines key open challenges, and provides a clear academic justification of how this thesis differs from prior works.
    
    \item \textbf{Chapter 3: Dataset Description and Formation} details the clinical origin and database architecture of MIMIC-IV version 3.1. It elaborates the multi-stage cohort filtering methodology, chunk-based extraction of ICU chart events, binary label assignment protocols, physiological plausibility filtering bounds, 24-hour observation windowing, exploratory data analysis of vital distributions, and explicit identification of dataset limitations.
    
    \item \textbf{Chapter 4: Proposed Methodology} formalizes the complete technical framework developed in this project. It provides mathematical formulations for ordinary least squares trend estimation, composite hemodynamic markers, standardized patient-level dataset partitioning, hyperparameter tuning grids, decentralized client allocation, PyTorch logistic regression parameterization, and step-by-step mathematical derivations of the Federated Averaging algorithm alongside complete procedural flowcharts and pseudocode.
    
    \item \textbf{Chapter 5: Results and Analysis} delivers an in-depth empirical evaluation of all trained models on the held-out test cohort. It provides detailed diagnostic breakdowns of accuracy, precision, recall, specificity, F1-score, and ROC-AUC, followed by an exhaustive analysis explaining why simple linear models outperformed complex tree ensembles on small clinical cohorts. Furthermore, it analyzes the collapse of isolated hospital models, tracks federated learning convergence curves across communication rounds, investigates IID versus Non-IID performance, reviews tree feature importances, and documents the operational validation of the deployed web application.
    
    \item \textbf{Chapter 6: Conclusion and Future Scope} summarizes the principal technical and clinical conclusions established by this research, outlines real-world system limitations, and proposes promising directions for future exploration, including differential privacy integration, multimodal waveform fusion, and clinical prospective deployment.
\end{itemize}
